\documentclass[10pt]{article}
\usepackage[T1]{fontenc}
\usepackage[utf8]{inputenc}
\usepackage[french]{babel}
\usepackage{amsmath,amssymb,mathrsfs,mathpazo}
\usepackage{geometry}
\geometry{a4paper,margin=2cm}
\usepackage{graphicx,array,tabularx,tabularray,colortbl,multirow,multicol,booktabs,enumitem,pifont,ulem,fancybox,diagbox,cancel}
\usepackage{pstricks,pstricks-add,pst-plot,pst-tree,pst-node,pst-text,pst-func,pst-math,pst-bspline,pst-3dplot,pst-eucl,pst-solides3d}
\let\pscalculate\undefined
\usepackage{pst-bezier}
\usepackage{tikz,pgfplots,tkz-tab}
\usetikzlibrary{arrows,arrows.meta,patterns,calc,positioning,trees,decorations.markings,shapes.geometric,shapes.misc,shapes.arrows,intersections,decorations.pathmorphing,decorations.pathreplacing,backgrounds,shadows,babel}
\pgfplotsset{compat=1.18}
\usepgfplotslibrary{fillbetween}
\usepackage[np]{numprint}
\usepackage{listings,makecell,pas-tableur,twoopt}
\let\barre\undefined
\usepackage[pstricks]{tablvar}
\definecolor{darkgreen}{RGB}{0,100,0}
\providecommand{\red}{\color{red}}
\providecommand{\blue}{\color{blue}}
\providecommand{\green}{\color{green}}
\providecommand{\black}{\color{black}}
\providecommand{\white}{\color{white}}
\providecommand{\gray}{\color{gray}}
\providecommand{\cyan}{\color{cyan}}
\providecommand{\magenta}{\color{magenta}}
\providecommand{\yellow}{\color{yellow}}
\providecommand{\violet}{\color{violet}}
\providecommand{\orange}{\color{orange}}
\providecommand{\euro}{\texteuro}
\providecommand{\EUR}{\texteuro}
\providecommand{\ieme}{\textsuperscript{e}}
\providecommand{\degres}{\ensuremath{^\circ}}
\providecommand{\degree}{\ensuremath{^\circ}}
\providecommand{\vv}[1]{\overrightarrow{#1}}
\providecommand{\Coord}[1]{#1}
\providecommand{\decosix}{}
\providecommand{\starredbullet}{\ensuremath{\bullet}}
\pagestyle{empty}
\setlength{\parindent}{0pt}
\usepackage{adjustbox}\DeclareUnicodeCharacter{2236}{:}\DeclareUnicodeCharacter{2190}{\ensuremath{\leftarrow}}
\providecommand{\R}{\mathbb{R}}
\providecommand{\N}{\mathbb{N}}
\providecommand{\Z}{\mathbb{Z}}
\providecommand{\Q}{\mathbb{Q}}
\providecommand{\D}{\mathbb{D}}
\providecommand{\C}{\mathbb{C}}
\providecommand{\e}{{\mathrm{e}}}
\providecommand{\vect}[1]{\overrightarrow{#1}}
\providecommand{\vectt}[1]{\overrightarrow{\mathrm{#1}}}
\providecommand{\vvt}[1]{\overrightarrow{\mathrm{#1}}}
\providecommand{\np}[2][]{{#2~#1}}
\providecommand{\vv}[1]{{\overrightarrow{#1}}}
\renewcommand{\d}{{\mathrm{d}}}
\providecommand{\degres}{{^\circ}}
\providecommand{\degre}{{^\circ}}
\providecommand{\blue}{\color{blue}}
\providecommand{\red}{\color{red}}
\providecommand{\enskip}{\quad}
\providecommand{\llbracket}{\lbrack\!\lbrack}
\providecommand{\rrbracket}{\rbrack\!\rbrack}
\providecommand{\euro}{{\text{€}}}
\providecommand{\no}{{\text{n°}}}
\providecommand{\interval}[2]{{[#1;#2]}}
\providecommand{\intervFO}[2]{{[#1;#2[}}
\providecommand{\dotfill}{\ldots}

\providecommand{\ds}{\displaystyle}
\providecommand{\textding}[1]{#1}
\providecommand{\Oij}{\(\mathrm{O};\vec{\imath},\vec{\jmath}\)}
\providecommand{\Oijk}{\(\mathrm{O};\vec{\imath},\vec{\jmath},\vec{k}\)}
\providecommand{\Ouv}{\(\mathrm{O};\vec{u},\vec{v}\)}
\providecommand{\intervFF}[2]{\ensuremath{[#1;#2]}}
\renewcommand{\euro}{\ifmmode\text{\texteuro}\else\texteuro\fi}
\DeclareUnicodeCharacter{00B0}{\ifmmode^{\circ}\else\textdegree\fi}
\DeclareUnicodeCharacter{20AC}{\ifmmode\text{\texteuro}\else\texteuro\fi}
% Énoncés : Éducation nationale. Transcription et corrigés : APMEP, Denis Vergès et contributeurs. Découpage et classement : CoopMaths, Sébastien Lozano, Éric Elter et contributeurs. Adaptation de la présentation : Maths Explorer.
% CC BY-SA : https://coopmaths.fr/www/about/
% https://coopmaths.fr/alea/static/e3c/2020/tex/e3c_2020_00_sujet38_4.tex
\begin{document}
\noindent\textit{APMEP / CoopMaths. CC BY-SA. Présentation adaptée par Maths Explorer.}\par\medskip
\medskip

\begin{enumerate}
\item Soit $f$ la fonction définie sur l'intervalle [0~;~2] par $f(x) = 8x - 2x^3$.
\begin{enumerate}
\item Montrer que pour tout réel $x$ de [0; 2], $f'(x)$ a le même signe que $4 - 3x^2$.
\item Étudier les variations de la fonction $f$ sur [0~;~2].
\end{enumerate}
\item Dans un repère orthonormal, on considère la parabole $p$ d'équation $y = x^2$ et la droite $\mathscr{D}$ d'équation $y = 4$.

On considère le rectangle $MSFE$ tel que :
\begin{itemize}
\item $M$ un point de $p$ dont l'abscisse $x$ est un réel de $]0~;~2[$.
\item $S$ est le symétrique de $M$ par rapport à l'axe des ordonnées.
\item $E$ et $F$ sont respectivement les projetés orthogonaux de $M$ et $S$ sur la droite $\mathscr{D}$.
\end{itemize}

\begin{center}



\psset{xunit=1cm,yunit=1cm,labelFontSize=\scriptstyle,showorigin=false}
\begin{pspicture}(-4.4,-1.5)(4.5,5.6)
\psframe[fillstyle=solid, fillcolor= blue](-1,1)(1,4)
\multido{\n=-4+0.2}{40}{\psline[linewidth=0.25pt,linecolor=lightgray](\n,-0.99)(\n,5.2)}
\multido{\n=-0.8+0.2}{30}{\psline[linewidth=0.25pt,linecolor=lightgray](-4,\n)(4,\n)}
\multido{\n=-4+1}{9}{\psline[linewidth=0.45pt](\n,-0.9)(\n,5.2)}
\multido{\n=0+1}{6}{\psline[linewidth=0.45pt](-4,\n)(4,\n)}
\psaxes[linewidth=0.95pt]{-}(0,0)(-4,-0.99)(4,5.2)
\def\Func{x x  mul  }
\psplot[plotpoints=2000,linewidth=0.85pt,linecolor=red]{-2.27}{2.27}{\Func}
\uput[dl](0,0){O}
\psdots[dotstyle=Bullet,dotscale=1.5,](1,1)(-1,1)(-1,4)(1,4)
\uput[dr](1,1){\footnotesize $M$}\uput[ur](1,4){\footnotesize $E$}\uput[ul](-1,4){\footnotesize $F$}\uput[dl](-1,1){\footnotesize $S$}
\psline[linewidth=1.25pt](-4,4)(4,4)\uput[u](3.8,4){$\mathscr{D}$}
\end{pspicture}
\end{center}

	\begin{enumerate}
		\item  Lorsque l'abscisse $x$ du point $M$ varie dans ]0~;~2[, l'aire du rectangle $MSFE$ est-elle constante ?
		\item Montrer que l'aire du rectangle $MSFE$ en fonction de l'abscisse $x$ de $M$ est $8x - 2x^3$.
		\item Montrer que l'aire maximale du rectangle $MSFE$ est $\dfrac{32}{3\sqrt{3}}$.
	\end{enumerate}
\end{enumerate}

\end{document}
